\documentclass[conference]{IEEEtran}

\usepackage{cite}
\usepackage{amsmath}
\usepackage{booktabs}
\usepackage{array}
\usepackage{url}

\title{EpiScope: A Big Data Analytics Pipeline for Synthetic Disease Surveillance, Outbreak Detection, and Preparedness Support}

\author{
\IEEEauthorblockN{Md. Samin Osman (0222220005101042) \\ Robiul Hassan (0222220005101043) \\ Walid Talal (0222220005101045) \\ Md. Mahafujul Alam (0222220005101066)}
\IEEEauthorblockA{
Department of Computer Science and Engineering\\
Premier University, Chattogram, Bangladesh
}
}

\begin{document}

\maketitle

\begin{abstract}
Public-health analytics can require the processing of large encounter datasets, data-quality controls, trend analysis, and the preparation of information for operational planning. This paper presents EpiScope, a prototype big-data pipeline that demonstrates these tasks using synthetic disease-encounter data. The architecture uses Apache Spark for data transformations, Bronze, Silver, and Gold processing layers, Parquet analytical storage, baseline outbreak detection and forecasting, resource-allocation and preparedness-calendar modules, and a D3.js dashboard export path. In a documented 25-million-row run, 24,825,210 records were retained in the clean Silver output and 174,790 were quarantined, reconciling to the input count. Documented processing-stage timings total 27.40 minutes, excluding synthetic-data generation. The latest saved detection evaluation reports 190 true positives, 25,336 false positives, 169 false negatives, and 554,145 true negatives, corresponding to 52.92\% sensitivity, 4.37\% false-positive rate, and 0.74\% precision. The latest saved 28-day forecast evaluation reports a mean absolute error of 2.0182 and a root mean squared error of 2.6986 across 8,960 district-disease-day observations. Results are limited to the synthetic dataset and current evaluation procedures. Historical metrics differ substantially from the latest saved evaluations; their original run configurations are unavailable. EpiScope is therefore presented as an educational and analytical prototype, not a validated public-health decision system.
\end{abstract}

\begin{IEEEkeywords}
big data analytics, Apache Spark, data quality, disease surveillance, outbreak detection, forecasting, privacy-aware visualization
\end{IEEEkeywords}

\section{Introduction}

Disease-surveillance workflows may involve large volumes of records from different sources and require reliable preparation before analytical results can be presented. A useful analytics prototype should demonstrate data ingestion, validation, aggregation, analytical processing, and the controlled publication of derived information.

EpiScope was developed as a big-data analytics project to demonstrate such a workflow on synthetic encounter records. It combines layered data processing with baseline detection and forecasting, resource-allocation and preparedness-calendar outputs, and a browser dashboard. The project emphasizes an end-to-end data path and records limitations in the evaluation and privacy checks.

The main contributions of this project are:
\begin{itemize}
    \item A Bronze, Silver, and Gold processing workflow for synthetic encounter data.
    \item A documented large-input data-quality reconciliation and quarantine process.
    \item Baseline outbreak-detection and fixed-horizon forecasting evaluations.
    \item Experiments documenting pipeline timing, Gold-stage scaling, and Spark configuration comparisons.
    \item A public-data preparation step that suppresses tested small case counts before dashboard export.
\end{itemize}

The paper does not claim clinical validity, real-world epidemiological accuracy, or formal privacy guarantees.

\section{System Architecture}

EpiScope separates data preparation, analytics, and dashboard serving. Synthetic encounter records are ingested into the Bronze layer. The Silver stage validates records and separates accepted rows from quarantined rows. The Gold stage creates aggregated analytical datasets. Detection, forecasting, resource allocation, and preparedness-calendar modules consume analytical outputs. Public-data preparation is then applied before JSON datasets are exported to the D3.js dashboard.

Apache Spark is used for data transformations, Parquet for processed analytical storage, and D3.js v7 for browser visualization \cite{spark,parquet,d3}. The project runs in a Google Colab environment with Spark in local mode for the recorded experiments. Consequently, measured runtimes depend on the runtime allocation, storage, partitioning, caching, and Spark configuration.

\section{Data and Processing Method}

\subsection{Synthetic Input}

The encounter records and outbreak labels used in the recorded experiments are synthetic. The project includes a generator for creating encounter records at different scales. Weather-related fields, if present in a given generated schema, must likewise be treated as synthetic or derived unless their provenance is independently established. Results should not be interpreted as evidence of performance on hospital records, government surveillance data, or real weather observations.

\subsection{Layered Processing}

The pipeline uses three logical data layers:
\begin{itemize}
    \item \textbf{Bronze:} ingestion of source encounter records.
    \item \textbf{Silver:} validation, cleaning, and quarantine of records that fail implemented checks.
    \item \textbf{Gold:} aggregated datasets for downstream analytics and serving.
\end{itemize}

The Silver stage aims to preserve accounting between accepted and quarantined rows. Reconciliation demonstrates count consistency but does not independently prove that every validation rule is correct.

\subsection{Incremental Inventory Input}

The project includes a separate synthetic daily inventory feed with facility stock and bed-occupancy fields. The Colab workflow generates a one-day snapshot and processes it through Bronze, Silver, and Gold. Bronze appends only unseen \texttt{(snapshot\_date, facility\_id)} keys, so replaying the same snapshot should append zero rows. The first-write-wins key policy does not apply corrections to a previously ingested facility-day. This incremental path is currently a parallel demonstration; the resource-allocation optimizer does not consume its Gold output. The submitted Colab notebook records a successful two-pass execution of the inventory pipeline.

\subsection{Analytics}

The outbreak detector is a prototype evaluated at the district-disease-day level for observations with at least 14 baseline days. Forecasting uses a seven-day-average baseline computed before the test period and repeated as a fixed prediction across a 28-day horizon. This is not a rolling forecast that updates predictions every day.

Resource allocation and preparedness-calendar components are included as analytical demonstrations; their outputs are not validated operational guidance.

\subsection{Public Data and Privacy}

A public-data preparation step precedes dashboard JSON export. The documented policy suppresses exact case counts from one through four. A recorded check of the public case export found no exposed values in that range in the tested output.

This is a limited check, not a complete privacy audit: inference may still be possible through geographic and temporal combinations, totals, linked views, or repeated queries.

\section{Experimental Setup}

The recorded experiments were run in Google Colab using Spark in local mode. The documented full-scale processing input contained 25,000,000 rows. Synthetic-data generation is excluded from the reported 27.40-minute processing-stage total.

Forecast evaluation used the test period from December 4 through December 31, 2025, with 8,960 district-disease-day observations. Detection metrics are taken from the latest saved evaluation output. Timing, optimization, scaling, and data-quality counts are documented historical measurements; raw timing logs were not available in the audited paths.

\section{Results}

\subsection{Data Quality}

Table~\ref{tab:quality} summarizes the documented reconciliation for the 25-million-row run.

\begin{table}[htbp]
\caption{Documented Silver-Stage Reconciliation}
\label{tab:quality}
\centering
\begin{tabular}{lr}
\toprule
Measure & Rows \\
\midrule
Input & 25,000,000 \\
Clean output & 24,825,210 \\
Quarantined output & 174,790 \\
Clean + quarantined & 25,000,000 \\
Quarantine rate & 0.69916\% \\
\bottomrule
\end{tabular}
\end{table}

The clean and quarantined counts sum to the input count. These values were transcribed from the project's documentation and were not independently recounted during report preparation.

\subsection{Pipeline Runtime}

Table~\ref{tab:runtime} presents documented stage runtimes. The total is a sum of recorded stages, not an independent timing measurement.

\begin{table}[htbp]
\caption{Documented Processing-Stage Runtime}
\label{tab:runtime}
\centering
\begin{tabular}{lr}
\toprule
Stage & Runtime (min) \\
\midrule
Bronze & 5.16 \\
Silver & 12.83 \\
Gold & 3.57 \\
Outbreak detection & 1.47 \\
Forecast generation & 0.94 \\
Resource allocation & 0.59 \\
Preparedness calendar & 0.57 \\
Public-data build & 0.50 \\
Dashboard export & 0.66 \\
Forecast evaluation & 0.86 \\
Detection evaluation & 0.75 \\
\midrule
Total measured stages & 27.40 \\
\bottomrule
\end{tabular}
\end{table}

The documented total is below the project's 45-minute target. Synthetic-data generation, recorded separately at approximately 7.17 minutes, is excluded. These measurements are environment-specific, and the raw timing logs were not found in the audited paths.

\subsection{Outbreak Detection}

Table~\ref{tab:detection} shows the latest saved detection evaluation.

\begin{table}[htbp]
\caption{Latest Saved Outbreak-Detection Evaluation}
\label{tab:detection}
\centering
\begin{tabular}{lr}
\toprule
Metric & Value \\
\midrule
True positives (TP) & 190 \\
False positives (FP) & 25,336 \\
False negatives (FN) & 169 \\
True negatives (TN) & 554,145 \\
Sensitivity & 52.92\% \\
False-positive rate & 4.37\% \\
Precision & 0.74\% \\
\bottomrule
\end{tabular}
\end{table}

Sensitivity indicates that the detector identifies approximately half of the positive evaluation cases. However, precision of 0.74\% means that false-positive alerts greatly outnumber true-positive alerts. This substantially limits the usefulness of the current detector without review and further tuning. It must not be used for unreviewed operational decisions.

\subsection{Forecasting}

The latest saved forecast evaluation contains 8,960 observations for December 4 to 31, 2025. The seven-day-average fixed-horizon baseline achieved MAE \(=2.0182\) cases and RMSE \(=2.6986\) cases. MAE expresses the average absolute difference between predicted and actual counts; RMSE weights larger errors more heavily. These values describe only the current synthetic evaluation and depend on the evaluation target and calendar-completion procedure.

Earlier project documentation reported MAE approximately 5.9894 and RMSE approximately 11.6049. Their source run and configuration are unavailable for independent verification, so these figures are identified as historical. The latest saved evaluation is reported separately as the current result.

\subsection{Optimization and Scaling}

Documented Gold-stage experiments report 4.91 minutes for AQE disabled with 200 shuffle partitions, compared with 3.45 minutes for AQE enabled with 32 partitions. The latter is approximately 29.7\% lower than the 4.91-minute baseline. Because these are historical single-run measurements without raw logs in the audited paths, they should be interpreted as indicative observations rather than a reliable general speedup.

Gold-stage runtimes were documented as 1.55 minutes for 6.25 million rows, 2.03 minutes for 12.5 million rows, and 3.45 minutes for 24,825,210 clean rows. The smaller datasets were generated using prefix or row-limit samples and may not preserve the full dataset distribution. These measurements describe the Gold stage only, not end-to-end pipeline scaling.

\subsection{Dashboard Export Snapshot}

The saved dashboard files contain 36,000 weekly case rows, 17,063 alerts, 8,960 forecasts, 64 district allocations totaling 100 units, and 3,840 preparedness-calendar rows. The combined JSON payload is 9.81 MiB. The weekly case counts range from 32 to 3,667, with no values from 1 through 4 in this saved file. These counts were read from the local JSON exports. The cases and alerts files have later modification times than the forecast, resource, and calendar files, so the payload is a mixed saved snapshot rather than a synchronized pipeline export.

\section{Discussion}

The project demonstrates an integrated workflow that connects large-input processing, data-quality reconciliation, baseline analytics, and dashboard-oriented export. The documented Silver-stage reconciliation provides evidence of row accounting for the recorded run. The runtime records suggest that the pipeline stages fit within the stated processing target in the measured environment, although the absence of raw timing logs limits independent verification.

The analytical results also show why evaluation metrics must be interpreted together. A sensitivity of 52.92\% may appear useful in isolation, but precision of 0.74\% indicates a high false-alert burden. Similarly, the forecast errors cannot establish real-world predictive utility because the evaluation data are synthetic and the baseline is intentionally simple. The difference between current and historical forecast metrics highlights the importance of experiment provenance and automated result capture.

\section{Limitations and Future Work}

The principal limitations are:
\begin{itemize}
    \item \textbf{Synthetic data:} Results do not establish performance on real surveillance data.
    \item \textbf{Detection quality:} Very low precision and a substantial false-positive count require investigation of label construction, thresholds, class imbalance, and evaluation logic.
    \item \textbf{Forecast baseline:} The fixed seven-day average does not adapt through the 28-day horizon and should be compared against additional baselines using a reproducible evaluation protocol.
    \item \textbf{Historical discrepancies:} Older forecast and detection figures differ from current saved outputs and have not been reconciled.
    \item \textbf{Benchmark evidence:} Historical runtime and optimization records lack raw logs in the audited paths, and the reported values may be affected by caching and environment variability.
    \item \textbf{Privacy:} Suppressing small counts is not a formal privacy guarantee. All public exports and cross-view inference risks require further testing.
    \item \textbf{Incremental inventory:} The synthetic daily feed demonstrates key-based replay protection, but the optimizer does not consume its outputs and corrected snapshots are not merged.
    \item \textbf{Quality-rule evidence:} The distinct implemented Silver validation rules and their test evidence can be compared with the assignment rule count to document implementation coverage.
\end{itemize}

Future work should include reproducible experiment logging, repeated benchmark runs, a source-derived data-quality rule inventory, expanded public-output privacy tests, alternative forecasting baselines, and detector threshold analysis. Any application to real public-health decisions would require representative real data, domain-expert review, and appropriate governance.

\section{Conclusion}

EpiScope demonstrates a layered big-data pipeline for synthetic disease-surveillance analytics, including data cleaning and quarantine, aggregation, baseline detection and forecasting, and dashboard data preparation. The documented large-input run reconciles 25 million input rows and reports 27.40 minutes across measured processing and evaluation stages, excluding data generation. The latest saved evaluations report 52.92\% detection sensitivity but only 0.74\% precision, and forecast MAE/RMSE of 2.0182/2.6986 on a synthetic 28-day test period.

These results are preliminary and have limitations in historical comparability and independent reproducibility. The system should be regarded as an educational prototype rather than a validated public-health decision-support tool.

\begin{thebibliography}{00}

\bibitem{spark}
Apache Software Foundation, ``Apache Spark Documentation.''
[Online]. Available: \url{https://spark.apache.org/docs/latest/}

\bibitem{parquet}
Apache Software Foundation, ``Apache Parquet Documentation.''
[Online]. Available: \url{https://parquet.apache.org/docs/}

\bibitem{d3}
M. Bostock, V. Ogievetsky, and J. Heer, ``D$^3$: Data-Driven Documents,''
\emph{IEEE Transactions on Visualization and Computer Graphics},
vol. 17, no. 12, pp. 2301-2309, Dec. 2011,
doi: 10.1109/TVCG.2011.185.

\bibitem{ieee}
IEEE, ``IEEE Reference Guide.'' [Online]. Available:
\url{https://journals.ieeeauthorcenter.ieee.org/your-role-in-article-production/ieee-editorial-style-manual/}

\end{thebibliography}

\appendices
\documentclass[11pt,a4paper]{article}
\usepackage[margin=1in]{geometry}

\begin{document}

\clearpage
\section*{Policy Brief for Public-Health Decision Makers}
\subsection*{EpiScope: A research prototype for surveillance and preparedness}

\noindent\textbf{Audience:} District health officials, surveillance teams, and emergency-preparedness planners.

\noindent\textbf{Prepared by:} Md. Samin Osman (0222220005101042), Robiul Hassan (0222220005101043), Walid Talal (0222220005101045), and Md. Mahafujul Alam (0222220005101066).

\subsection*{The decision problem}
Health teams need to notice unusual disease activity early enough to investigate it and prepare staff, supplies, and services. Relevant information may come from different systems and arrive late or with inconsistent quality. A useful analytics workflow must therefore check incoming data, summarize trends, provide interpretable signals, and communicate limitations alongside recommendations.

EpiScope demonstrates this workflow using synthetic encounter and facility data. It places records into Bronze, Silver, and Gold processing stages; produces outbreak signals and short-term forecasts; and prepares resource and calendar outputs for a D3.js dashboard. A separate demonstration generates a daily synthetic stock and bed-occupancy snapshot and processes it through Bronze, Silver, and Gold with replay protection. The snapshot is synthetic facility information and the resource-allocation optimizer does not consume it.

\subsection*{What the current evidence says}
The saved 25-million-row synthetic run records 24,825,210 accepted records and 174,790 quarantined records, which reconcile to the input. Documented processing-stage timings total 27.90 minutes, excluding data generation. These are historical Google Colab measurements; raw timing logs are unavailable for independent reproduction.

The saved synthetic detection evaluation reports 52.92\% sensitivity, a 4.37\% false-positive rate, and 0.74\% precision. Sensitivity is the share of labeled outbreak cases detected; the false-positive rate is the share of non-outbreak observations incorrectly flagged; precision is the share of alerts that match a labeled outbreak. Thus, most alerts in this evaluation were false positives and would require substantial human review. Alerts must not be sent to field teams as confirmed events.

The saved forecasting evaluation reports a mean absolute error of 2.0182 cases and a root mean squared error of 2.6986 cases across 8,960 district-disease-day observations from December 4 to 31, 2025. These measures summarize the difference between predicted and recorded case counts; root mean squared error gives larger misses more weight. The forecast repeats a seven-day average through a fixed 28-day horizon and has not been shown to outperform a seasonal baseline. All figures are historical outputs from synthetic data, not evidence of real-world outbreak performance or health impact.

\clearpage
\subsection*{Recommended use and safeguards}
\begin{enumerate}
    \item \textbf{Keep a person in the decision loop.} Treat an alert as a prompt to review recent counts, reporting completeness, local context, and possible data errors. Do not treat it as confirmation of an outbreak.
    \item \textbf{Improve alert quality before expanding use.} Test alternative thresholds on a held-out period. Review sensitivity, precision, false-positive rate, and alert volume together, and agree with district teams on an acceptable review workload.
    \item \textbf{Compare forecasts with simple alternatives.} Evaluate the seven-day-average method against a seasonal-naive baseline and report errors by disease, district, and time period. Identify where forecasts are unreliable.
    \item \textbf{Check practical constraints before allocating resources.} Validate stock, bed capacity, staff, transport, and local protocols. Current allocation and preparedness outputs are demonstrations based on illustrative assumptions, not validated operational instructions.
    \item \textbf{Protect data and review disclosure risks.} Restrict access to encounter-level data. The Silver processing stage hashes patient IDs before producing clean records; identifiers in Bronze and quarantine data remain sensitive and require access controls. Public case counts from one through four are suppressed. Suppression reduces direct exposure but does not guarantee anonymity; linked totals and views may still reveal small counts.
\end{enumerate}

\subsection*{Conditions before operational use}
The data in EpiScope are synthetic. Encounter labels, facility snapshots, and temperature and rainfall values are generated rather than collected from health facilities or verified weather stations. Synthetic data are useful for testing the processing path, but cannot establish how the system performs under real reporting delays, missing facility submissions, changing disease patterns, or local resource constraints.

Before any field use, the system would need representative and authorized data, independent epidemiological review, stronger detector precision, forecast comparison and calibration, a tested privacy and access-control plan, and a defined process for alert review and escalation. The incremental inventory demonstration would also need validated live sources, a policy for corrected or late snapshots, and integration with allocation logic before it could inform supply recommendations.

\vspace{1.5em}
\noindent\textbf{Decision requested:} Continue supervised validation and improve the evidence. Do not use EpiScope for autonomous outbreak declarations, clinical decisions, or unreviewed resource allocation.

\clearpage

\end{document}


\end{document}
